\section{Star Differeth from Star}

In the chapter on the resurrection, the Apostle looks up at the night sky and reads in it a law of glory: ``One is the glory of the sun, another the glory of the moon, and another the glory of the stars. For star differeth from star in glory. So also is the resurrection of the dead.''\endnote{1 Cor 15:41--42, Douay--Rheims (Challoner).} The Church teaches the same inequality of glory. The Council of Florence, in the sentence that confesses the vision of God himself, three and one, as he is, defines that the blessed see him \latin{pro meritorum tamen diversitate alium alio perfectius}: yet according to the diversity of merits, one more perfectly than another.\endnote{Council of Florence, \emph{Laetentur caeli} (1439), Denzinger n.~1305.} The Council of Trent, in its decree on justification, anathematizes anyone who denies that the justified man, by good works done through the grace of God and the merit of Jesus Christ, truly merits an increase of grace, eternal life, and \latin{gloriae augmentum}, an increase of glory.\endnote{Council of Trent, session VI (13 January 1547), decree on justification, canon 32, in Denzinger, \emph{Enchiridion symbolorum}, n.~1582 (Latin).} The same canon calls these works both the gifts of God and the merits of the man, so the differences in glory are differences in grace received and used.

\subsection{What differs, and what does not}

Thomas Aquinas is clear that the object of the vision is the same for all: ``the same object will be presented to all---viz.\ the essence of God.''\endnote{Aquinas, \emph{Summa theologiae} I, q.~12, a.~6, reply to obj.~3, trans.\ English Dominican Province.} No one in heaven has a lesser God, or sees him through a dimmer likeness, or is given a part where another has the whole. What differs is the seeing: the strength of the light of glory in each, by which the one God is seen more or less perfectly. Nor is that strength measured by natural gifts, as though the clever were to see more than the simple; Thomas raises the objection himself and answers that the diversity arises ``on the part of the diverse faculty of the intellect, not, indeed, the natural faculty, but the glorified faculty.''\endnote{Aquinas, \emph{Summa theologiae} I, q.~12, a.~6, obj.~3 and reply, trans.\ English Dominican Province.}

The Council names merit as the measure. Thomas explains the measure by charity:

\begin{quote}
Hence the intellect which has more of the light of glory will see God the more perfectly; and he will have a fuller participation of the light of glory who has more charity; because where there is the greater charity, there is the more desire; and desire in a certain degree makes the one desiring apt and prepared to receive the object desired. Hence he who possesses the more charity, will see God the more perfectly, and will be the more beatified.\endnote{Aquinas, \emph{Summa theologiae} I, q.~12, a.~6, trans.\ English Dominican Province.}
\end{quote}

Merit, on this explanation, is not credit laid up to be exchanged at the end for a better seat: a life of charity enlarges the soul, and the enlargement is what shows in glory. Cyprian had tied the reward to desire a thousand years before, at the end of his treatise on the plague, praying that Christ would look upon the eager longing of his people, ``He who will give the larger rewards of His glory to those whose desires in respect of Himself were greater!''\endnote{Cyprian, \emph{On the Mortality} 26, trans.\ Robert Ernest Wallis, \emph{The Ante-Nicene Fathers}, vol.~5 (Buffalo: Christian Literature Publishing Company, 1886).} Charity is poured into the heart by God and hidden from every eye but his. No one can measure his own, still less his neighbour's, so the doctrine of degrees gives no ground for ranking anyone on earth; the rank of glory is known only in heaven.

\subsection{Degrees without envy}

Even a child can raise the objection that an unequal heaven must leave someone sad. Augustine answered it in the last chapter of \emph{The City of God}:

\begin{quote}
But who can conceive, not to say describe, what degrees of honour and glory shall be awarded to the various degrees of merit? Yet it cannot be doubted that there shall be degrees. And in that blessed city there shall be this great blessing, that no inferior shall envy any superior, as now the archangels are not envied by the angels, because no one will wish to be what he has not received, though bound in strictest concord with him who has received; as in the body the finger does not seek to be the eye, though both members are harmoniously included in the complete structure of the body. And thus, along with his gift, greater or less, each shall receive this further gift of contentment to desire no more than he has.\endnote{Augustine, \emph{The City of God} XXII.30, trans.\ Dods (1871).}
\end{quote}

Thomas says the same of joy: each of the blessed reaches the term appointed for him, and ``each one's joy will be full with regard to himself, because his desire will be fully set at rest; yet one's joy will be greater than another's, on account of a fuller participation of the Divine happiness.''\endnote{Aquinas, \emph{Summa theologiae} II-II, q.~28, a.~3, reply to obj.~2, trans.\ English Dominican Province.} Thérèse of Lisieux learned it as a child from her sister Pauline, and told it in her own story:

\begin{quote}
One day I expressed surprise that God does not give an equal amount of glory to all the elect in Heaven---I was afraid that they would not all be quite happy. She sent me to fetch Papa's big tumbler, and put it beside my tiny thimble, then, filling both with water, she asked me which seemed the fuller. I replied that one was as full as the other---it was impossible to pour more water into either of them, for they could not hold it.\endnote{Thérèse of Lisieux, \emph{The Story of a Soul}, ch.~2, in \emph{Soeur Thérèse of Lisieux}, ed.\ T.~N. Taylor (London: Burns, Oates \& Washbourne, 1912).}
\end{quote}

In ``the slowest sphere'' of the \emph{Paradiso} Dante meets Piccarda, torn from her cloister by violent men, and asks whether the souls placed so low do not desire a higher place. She smiles, and answers:

\begin{quote}
Brother, our will is quieted by virtue\\
\quad Of charity, that makes us wish alone\\
\quad For what we have, nor gives us thirst for more.\\
\ldots\\
And his will is our peace; this is the sea\\
\quad To which is moving onward whatsoever\\
\quad It doth create, and all that nature makes.\endnote{Dante Alighieri, \emph{Paradiso} III.51, 70--72, 85--87 (the cloister, III.106--107), trans.\ Henry Wadsworth Longfellow (1867), Project Gutenberg eBook 1004.}
\end{quote}

And the poet understands: ``Then it was clear to me how everywhere / In heaven is Paradise, although the grace / Of good supreme there rain not in one measure.''\endnote{Dante, \emph{Paradiso} III.88--90, trans.\ Longfellow.} Gregory the Great, speaking of the light in which the blessed know one another, uses a phrase that bears on the fear of envy too: they ``do in that place with unspeakable brightness (common to all) behold God.''\endnote{Gregory the Great, \emph{Dialogues} IV.33, in \emph{The Dialogues of Saint Gregory}, trans.\ P.~W.\ (Paris, 1608), ed.\ Edmund G. Gardner (London: Philip Lee Warner, 1911).} If the brightness is common to all, no one possesses God at another's expense.

The doctrine of degrees therefore asks something of the Christian and also consoles him. Every act of charity done now widens the soul for what it will receive. The measure is love, which is within the reach of the least, and every vessel in heaven will be full.
